\begin{aidisclaimer} * % ← Remove this asterisk, if you have not used AI in creating this thesis.

    I hereby declare, that the AI-based applications used in generating this work are as follows:

    \begin{center}
        \begin{tabularx}{\linewidth}{X|l}
            \toprule
            \textbf{Application}                                                     & \textbf{Version}                                               \\
            \midrule
            GitHub Copilot chat \parencite{githubAskingGitHubCopilot}                & Claude Sonnet 4.5 \parencite{anthropicIntroducingClaudeSonnet} \\
            GitHub Copilot code suggestions \parencite{githubGettingCodeSuggestions} & Grok Code Fast 1 \parencite{xaiGrokCodeFast}                   \\
            Ollama local LLM  \parencite{Ollama}                                     & gpt-oss-20b \parencite{openaiGptoss20bModelOpenAI}             \\
            OpenAI API    \parencite{APIPlatformOpenAI}                              & gpt-5-nano \parencite{openaiGPT5NanoModel}                     \\
            Ollama local embeddings \parencite{Ollama}                               & nomic-embed-text-v1.5 \parencite{nomicUnboxingNomicEmbed}      \\
            \bottomrule
        \end{tabularx}
    \end{center}

    \section*{Purpose of the use of AI}

    I have used AI-based applications to assist in writing and reviewing code, generating LaTeX snippets, and improving the clarity and grammar of the text. The AI tools have been used as aids to enhance productivity and quality, but the core ideas, research, and conclusions of this thesis are entirely my own. This thesis and the subsequent RAG system was written in VS Code with GitHub Copilot chat and code suggestions enabled.

    Claude Sonnet 4.5 was primarily used to discuss, format, and structure text sections. The model gave suggestions for improving clarity, grammar, and coherence. It also helped in generating LaTeX code snippets for formatting. Grok Code Fast 1 was used in developing the resulting application code, providing suggestions and improvements to the code structure and logic in in-line editing. However, the software design, architecture, and implementation decisions were made by me based on my own understanding and research.

    Ollama local LLM (gpt-oss-20b) and OpenAI API (gpt-5-nano) were used as part of the RAG system developed in this thesis. They were employed to demonstrate the capabilities of retrieval-augmented generation in answering questions based on a custom knowledge base. Nomic-embed-text was used to generate text embeddings for the knowledge base documents. The design and implementation of the RAG system, including the choice of models, retrieval mechanisms, and evaluation methods, were my own. These models were not directly involved in writing the thesis text.

    \section*{Parts of this work, where AI was used}

    The AI tools were used in the following parts of this thesis:
    \begin{itemize}
        \item Advice on structuring and formatting all sections of the thesis.
        \item Reviewing and improving the clarity and grammar of all written text.
        \item Generating LaTeX code snippets for formatting tables and figures.
        \item Assisting in writing code for the RAG system implementation.
        \item Generating text embeddings for the knowledge base documents in the RAG system.
        \item Answering questions based on the knowledge base when using the RAG system.
    \end{itemize}

\end{aidisclaimer}
